संच~$A$ सांत असेल, तर त्याचे आकारमान नैसर्गिक संख्येने मोजता येते. अशा संख्यांची तुलना करून आकारमानांची तुलना करता येते. संच अनंत असतील, तर ही गोष्ट अधिक गुंतागुंतीची होते. अनंततेची पहिली पातळी म्हणजे \emph{गणनीय अनंत} संचांची पातळी. रिक्त संच गणनीय आहे. अरिक्त संच~$A$ च्या घटकांचे \emph{प्रगणन}, म्हणजे एकदिश अनंत यादी, मांडता येत असेल, अर्थात आच्छादक फलन~$f\colon \PosInt \to A$ अस्तित्वात असेल, तर तो संच गणनीय असतो. तो गणनीय असून सांत नसेल, तर गणनीय अनंत असतो. कँटर यांच्या \emph{नागमोडी पद्धतीने} गणनीय अनंत संचांच्या घटकांच्या जोड्यांचा संचही गणनीय आहे हे दिसते. त्याच्या आधारे परिमेय संख्यांचा संच~$\Rat$ देखील गणनीय आहे हे दाखवता येते.

काही अनंत संच गणनीय नसतात. अशा संचांना \emph{अगणनीय} म्हणतात. एखादा संच अगणनीय आहे हे थेट \emph{कर्णरेषीय युक्तिवादाने} किंवा \emph{न्यूनीकरणाने} दाखवता येते. कर्णरेषीय युक्तिवादात आपण विचाराधीन संच~$A$ गणनीय आहे असे गृहीत धरतो. त्याच्या एका काल्पनिक प्रगणनाच्या आधारे $A$ चा असा एक घटक परिभाषित करतो की व्याख्येमुळेच तो प्रगणनातील प्रत्येक घटकापेक्षा वेगळा ठरतो. त्यामुळे ते संपूर्ण $A$ चे प्रगणन असू शकत नाही आणि $A$ चे कोणतेही प्रगणन अस्तित्वात असू शकत नाही हे सिद्ध होते. न्यूनीकरणात $A$ च्या प्रत्येक घटकाची, अगणनीय असल्याचे माहीत असलेल्या एखाद्या संच~$B$ च्या घटकाशी, आच्छादक फलनाने संगती लावून $A$ अगणनीय असल्याचे दाखवतात. अशी संगती लावता आली, तर $A$ च्या काल्पनिक प्रगणनापासून $B$ चे प्रगणन मिळेल. $B$ अगणनीय असल्यामुळे $A$ चे प्रगणन अस्तित्वात असू शकत नाही.

सर्वसाधारणपणे अनंत संचांची आकारमानानुसार तुलना करता येते. $A$ आणि $B$ यांच्यात एकास-एक आच्छादन असेल, तर त्यांचे आकारमान समान असते, म्हणजेच ते \emph{तुल्यबल} असतात. $A$ पासून $B$ पर्यंत एकास-एक फलन असेल, तर $A$ आकारमानाने $B$ पेक्षा मोठा नाही ($\cardle{A}{B}$), अशीही व्याख्या करता येते. श्रेडर--बर्नस्टाइन प्रमेयानुसार $\cardle{A}{B}$ आणि $\cardle{B}{A}$ असतील, तर $A$ आणि $B$ तुल्यबल असतात. अशा रीतीने संचांच्या आकारमानांचा क्रम मिळतो. शेवटी, \emph{कँटरच्या प्रमेयानुसार} कोणत्याही $A$ साठी $\cardless{A}{\Pow{A}}$. $\Pow{\PosInt}$ अगणनीय आहे या निष्कर्षाचे हे सामान्यीकरण आहे. यातून अनंततेच्या केवळ दोनच पातळ्या नसून अनंत पातळ्या आहेत हे दिसते.
